\chapter{Writing a Radiology Report}

\label{ch:write_report}

The radiology report is the primary communication tool between the radiologist and the referring clinician. A high-quality report must be clear, concise, systematic, and clinically relevant. This chapter details the structured components of a professional radiology report, the use of standardized lexicons, and the protocols for communicating critical findings.

---

\section*{1. The Structure of a Radiology Report}
A standard, professional radiology report is divided into six distinct, structured sections:

\subsection*{Demographics and Exam Header}
Contains patient name, medical record number, date of birth, date of exam, accession number, and the specific modality/projection performed.

\subsection*{Clinical Indication}
A brief summary of the patient's symptoms, clinical signs, and the diagnostic question being asked (e.g., *"Right lower quadrant pain, rule out acute appendicitis"*). 

\subsection*{Technique and Protocol}
Detailed documentation of how the scan was performed. This has significant clinical and legal value:
\textbullet modality details (e.g., *"Helical CT of the abdomen and pelvis"*).
\textbullet administration of contrast (e.g., *"with 100 mL of Omnipaque-350 intravenous contrast during the portal venous phase"*).
\textbullet patient safety parameters (e.g., *"with oral contrast"* or *"using low-dose pediatric protocols"*).

\subsection*{Comparison}
List any prior imaging studies compared with the current scan. If no prior studies are available, state: *"None available."* Comparing current findings with past scans is critical for determining if a lesion is acute or stable.

\subsection*{Findings}
The main body of the report. It should be written using a \textbf{systematic, organ-based or anatomical sequence} (e.g., listing airways, lungs, pleura, heart, bones for chest X-ray).
\textbullet Describe the presence, size, location, and attenuation/signal characteristics of any abnormalities.
\textbullet Use precise anatomical terms and objective measurements (avoid vague descriptions like *"large"* or *"small"*).
\textbullet Document relevant negatives (e.g., stating *"no pneumothorax"* is crucial in a chest trauma scan).

\subsection*{Impression}
The most important section of the report. It is the synthesis of the findings into a clinical conclusion:
\textbullet Answer the clinical question asked in the indication section.
\textbullet List the primary diagnosis first, followed by key differentials in order of clinical likelihood.
\textbullet Keep it concise (normally no more than 3--5 bulleted points).
\textbullet Provide actionable recommendations for follow-up testing (e.g., biopsy, follow-up scan, or clinical correlation) when appropriate.

---

\section*{2. Standardized Reporting Lexicons}
To reduce ambiguity, modern radiology relies on standardized scoring systems and vocabularies (\textbf{RADS}):
\textbullet \textbf{BI-RADS} (Breast): Guide for mammography, breast ultrasound, and breast MRI.
\textbullet \textbf{LI-RADS} (Liver): Categorization of liver nodules in patients with cirrhosis/hepatitis B (assessing hepatocellular carcinoma risk).
\textbullet \textbf{Lung-RADS} (Lung): Standardized management of nodules detected during low-dose CT lung cancer screening.
\textbullet \textbf{PI-RADS} (Prostate): Grading of lesions on multiparametric prostate MRI.
\textbullet \textbf{TI-RADS} (Thyroid): Risk stratification of thyroid nodules on ultrasound.

---

\section*{3. Communication of Critical Findings}
Radiologists have a clinical duty to communicate life-threatening or unexpected findings directly to the referring physician in a timely manner.
\textbullet \textbf{Category 1 (Immediate - within minutes)}: Life-threatening findings requiring emergent intervention (e.g., tension pneumothorax, free air in the abdomen, acute aortic dissection, intracranial hemorrhage, misplaced endotracheal tube).
\textbullet \textbf{Category 2 (Urgent - within hours)}: Findings that are not immediately life-threatening but require prompt medical attention (e.g., appendicitis, testicular torsion, new pulmonary embolism, critical vascular stenosis).
\textbullet \textbf{Documentation}: The report must document that the critical finding was communicated, noting the name of the physician notified, the date, and the exact time of the call.

---

\section*{4. Standard Report Template}

This template is standard in clinical practice:

\begin{testbox}[RADIOLOGY REPORT TEMPLATE]
\begin{tabbing}
\textbf{Patient Name}: \= John Doe \hspace{4cm} \= \textbf{MRN}: 12345678 \\
\textbf{Exam Date}: \> August 1, 2026 \> \textbf{DOB}: 01/01/1980 \\
\textbf{Accession}: \> ACC99988 \> \textbf{Exam}: CT Abdomen \& Pelvis with IV Contrast
\end{tabbing}

\noindent\rule{\textwidth}{0.4pt}

\textbf{INDICATION}: \\
Abdominal pain. Rule out diverticulitis.

\textbf{TECHNIQUE}: \\
Axial CT images of the abdomen and pelvis were obtained following the administration of 100 mL of Omnipaque-350 intravenous contrast. Coronal and sagittal reformations were reviewed.

\textbf{COMPARISON}: \\
None.

\textbf{FINDINGS}: \\
\textbf{Lungs/Pleura}: Visualized lung bases are clear. No pleural effusion.\\
\textbf{Solid Organs}: Liver, spleen, pancreas, gallbladder, and adrenal glands are unremarkable. The kidneys demonstrate symmetric nephrograms without hydronephrosis.\\
\textbf{Gastrointestinal Tract}: The stomach and small bowel loops are normal in caliber. In the left lower quadrant, there is segmentally thickened sigmoid colon wall with multiple out-pouchings (diverticulum). Moderate surrounding mesenteric fat stranding is present. No fluid collection or abscess. The appendix is normal.\\
\textbf{Peritoneum \& Vessels}: No free fluid or free air in the peritoneal cavity. Abdominal aorta is normal in caliber.\\
\textbf{Bones \& Soft Tissues}: Degenerative changes of the lumbar spine.

\textbf{IMPRESSION}: \\
1. Acute, uncomplicated sigmoid diverticulitis. No evidence of bowel perforation or abscess.
\end{testbox}

\clearpage
